% RESEARCH_DRAFT_SHA256: 78010940e04b8d2333bd82742ecdc7ebd6934dc22fcb4768e5d92f0076e196ea
% Layout and source-reference wrappers expanded for the web renderer.
\documentclass[11pt]{article}
\usepackage[margin=25mm]{geometry}
\usepackage{fontspec,unicode-math}
\setmainfont{Latin Modern Roman}
\setmathfont{Latin Modern Math}
\usepackage{amsmath,amssymb,mathtools,xurl,hyperref}
\hypersetup{colorlinks=true,urlcolor=blue,linkcolor=blue}
\setcounter{secnumdepth}{0}
\setlength{\parindent}{0pt}
\setlength{\parskip}{5pt}
\setlength{\emergencystretch}{3em}
\newcommand{\sref}[2]{\hyperlink{source:#1:#2}{[S#2]}}
\newcommand{\cataloguescope}[1]{\paragraph{Source record.}#1}
\newcommand{\floor}[1]{\left\lfloor#1\right\rfloor}
\begin{document}
\section*{Research-attempt provenance}
This standalone research attempt retains the complete original definition body below,
including its historical source and preparation notes. Those notes describe the
earlier catalogue preparation; the new attempt and its verification are identified
separately after that body. No original ID, catalogue status, or source statement is changed.

\section{Exact ridge reduction for Lévy PIDEs}\label{statement}

\subsection{Definitions and mathematical statement}

\paragraph{Shared notation.}Standard number systems are $\mathbb N,\mathbb Z,\mathbb Q,\mathbb R,\mathbb C$. Specialized notation and index ranges follow the cited source and the definitions below.


A partial integro-differential equation combines derivatives of an unknown function with a nonlocal integral term. A Lévy measure $\nu$ describes jump sizes and satisfies $\int(1\wedge\|z\|^2)\,\nu(dz)<\infty$. A viscosity subsolution or supersolution satisfies the corresponding differential inequalities when touched by suitable test functions. Comparison bounds every subsolution by every supersolution with compatible boundary data, often yielding uniqueness.
\paragraph{Record scope and qualifications.}
The catalogue title is retained as metadata. Named certificates, contracts, bounds or frameworks in that title are not asserted as results of the original source extract.
This is a participant research agenda rather than a fixed PIDE. Dynamics, Lévy measures, boundary conditions, and the requested numerical or regularity guarantee are not specified. The supplied catalogue marks this record open; its present research status was not independently checked.
The source text below is rendered literally where necessary; damaged formulas and composite extracts are not silently repaired into a new mathematical claim.
\paragraph{Statement recorded in the supplied source.}
A.1 Arisawa, Mariko

We would like to introduce a viscosity solution's approach to a class of second-order partial-integro-differential equations (PIDE, in short) arising in the mathematical finance. The usual brownian models for the processes of the asset prices are replaced by the jump-diffusion processes. The derived PIDEs contain in the integral term, various Le'vy measures which correspond to the frequencies and the sizes of the jumps. In addition, according to the financial settings such as the American option, the look back option, the problem with transaction costs, a variety of nonlinear PIDEs are engendered. We shall introduce a new definition of the viscosity solution for such problems, and shall give the comparison, the existence, and the regularity of the solutions of viscosity solutions in this framework. As far as the high-dimensional numerical analysis, a very few is known at this stage, while from the practical view point a lot of important problems should be explored, such as the modelization of the high-dimensional Le'vy measures arising in the mathematical finances. This is the problem which I would like to discuss with the participants of the workshop. \href{https://aimath.org/WWN/optimalcontrol/optimalcontrol.pdf}{[S2]}
\paragraph{Source status.}
The supplied catalogue labels this source record open; that metadata was checked against the supplied snapshot on 2026-10-11. The label is preserved for provenance and is not an independent verification of present open status or mathematical correctness.
\subsection{Short English statement}
The participant proposes viscosity-solution theory and high-dimensional numerical methods for jump models in finance.
\subsection{Sources}
\begin{itemize}
\item[S1] Ulam AI and the UnsolvedMath Contributors, UnsolvedMath: A Curated Collection of Open Mathematics Problems, catalogue entry. Catalogue curation and metadata: CC BY 4.0. The source extract is retained; definitions, English summaries and qualifications are editorial additions.
\url{https://huggingface.co/datasets/ulamai/UnsolvedMath}.
\item[S2] AIM Workshop Problem Lists, source linked by the catalogue.
\url{https://aimath.org/WWN/optimalcontrol/optimalcontrol.pdf}.
\end{itemize}
\paragraph{Preparation and review.}
The original source and UnsolvedMath contributors retain attribution for the source extract. Codex (GPT-6; exact serving revision not exposed) assisted with definitions, English summaries, scope notes and formatting. No new human mathematical authorship or independently verified solution is asserted. Catalogue identity checked against the complete supplied registry and current repository. Distinct upstream records, including repeated questions and broad or damaged extracts, are retained for complete coverage of the source CS category. A repeated record is not a new independent conjecture. Complete mathematical-equivalence and current-literature review has not been performed. Preparation and final review follow contribution prompts C and D, with source gaps disclosed. Structural and rendering checks do not establish a new theorem or open status.
\label{end:statement}

\section{Literature and context check}
The AIM optimal-control workshop PDF, Arisawa's participant contribution,
was consulted. It is an agenda covering several nonlinear PIDEs, not one fixed
operator and boundary-value problem. The following constant-coefficient smooth
generator calculation tests a dimension-reduction route without claiming the
agenda's viscosity comparison, existence, regularity or numerical guarantees.
\section{Research attempt: Projection of jump generators needs a drift correction}
Choose an explicit auxiliary model on $\mathbb R^d$:
\[
 Lu=b\cdot\nabla u+\tfrac12\operatorname{Tr}(aD^2u)
 +\int\big[u(x+z)-u(x)-\mathbf1_{\|z\|\le1}\nabla u(x)\cdot z\big]\nu(dz),
\]
where $a\succeq0$ is constant and $\nu$ is a Levy measure. For a linear map
$P:\mathbb R^d\to\mathbb R^k$ take $u(x)=v(Px)$ with $v\in C_b^2$.
Push $\nu$ forward on the nonzero projected jumps, omitting any mass at $Pz=0$;
call the measure $\mu$. It is a Levy measure since
\[
 1\wedge\|Pz\|^2\le\max(1,\|P\|^2)(1\wedge\|z\|^2).
\]
The chain rule gives diffusion matrix $PaP^T$. To write the jump part using
the standard truncation $\mathbf1_{\|w\|\le1}$ in projected coordinates, one
must use drift
\[
 b'=Pb+\int\big(\mathbf1_{\|Pz\|\le1}-\mathbf1_{\|z\|\le1}\big)Pz\,\nu(dz).
\]
The integral is absolutely convergent. On $\|z\|>1$, only terms with
$\|Pz\|\le1$ contribute, bounded by one against finite measure. On
$\|z\|\le1$, disagreement requires $\|Pz\|>1$, hence
$\|z\|>1/\|P\|$ when $P\ne0$; this is a finite-measure annulus and the
integrand is bounded by $\|P\|$. The case $P=0$ contributes zero. Adding and
subtracting the projected truncation inside the absolutely defined compensated
integral now proves the formula. Small-jump integrability follows from the
quadratic Taylor bound for $C_b^2$ functions; large jumps use boundedness.

The proposed shortcut omitted the drift correction. Set $d=2,k=1$,
$P(z_1,z_2)=z_1$, $b=0$, $a=0$, and
$\nu=\lambda\delta_{(\varepsilon,2)}$ with $0<\varepsilon<1$.
The original jump is large, so the exact projected action is
$\lambda[v(y+\varepsilon)-v(y)]$. The projected jump is small, so using
$\mu=\lambda\delta_\varepsilon$ with zero drift instead produces
$\lambda[v(y+\varepsilon)-v(y)-\varepsilon v'(y)]$.
Their difference is $\lambda\varepsilon v'(y)$, nonzero for $v(y)=\sin y$
at $y=0$. The required corrected drift is precisely $\lambda\varepsilon$.

Even with correct truncation, a ridge ansatz need not close for state-dependent
coefficients. With no jumps or diffusion, $b(x)=(\sin x_2,0)$ and
$u(x)=\sin x_1$, one has $Lu=\sin x_2\cos x_1$. At the same projected
coordinate $x_1=0$, this equals zero at $x_2=0$ and one at $x_2=\pi/2$.
It is not a function of $Px$. Thus an exact lower-dimensional equation needs
coefficient compatibility on projection fibers, absent from the broad source.
\section{Verification and adversarial audit}
All compensation signs were checked by the one-atom example and a bounded
smooth test function. The pushforward integrability proof includes jumps in
the kernel and maps of arbitrary operator norm. The failure-of-closure example
uses bounded smooth coefficients, so unbounded drift is not its cause.
\paragraph{Audit corrections.} The compensator-dependent drift was retained and its sign tested; a smooth auxiliary generator identity was not presented as a viscosity theorem or a solution of the agenda.
\paragraph{External computation.} None; symbolic arguments and boundary checks only.
\section{Scope and remaining gaps}
Only a smooth constant-coefficient auxiliary identity was proved. The
source supplies no specific PIDE, controls, payoff or boundary conditions; no
viscosity comparison or high-dimensional numerical convergence result follows.
The first gap is a fully specified original model whose coefficients and data
support the proposed reduction.
\paragraph{Integration review.} Codex (GPT-6; exact serving revision not exposed) reviewed the written arguments and expanded layout and source-reference wrappers. Repository metadata and the production web renderer were checked using Python 3.13 and JavaScript with MathJax 3.2.2. This model-assisted review is not independent human proof verification; no PDF compilation or formal proof checking is asserted.

\section{Completion Estimate}
\noindent\textbf{COMPLETION: 0\%}
The calculations below the preserved source statement investigate a concrete route,
but establish no substantiated new progress toward the entire original target.
Known special cases and failed sufficient conditions are not credited as a solution.
\section{References}
The original source references and catalogue attribution are preserved in the
Sources section above. The following URLs were actually consulted for this attempt;
retention of an additional original reference does not claim fresh consultation of it.
\begin{itemize}
\item \url{https://aimath.org/WWN/optimalcontrol/optimalcontrol.pdf} (consulted 2026-10-11). 
\end{itemize}
\section{Model and contributor disclosure}
OpenAI gpt-6-astra, requested reasoning configuration ultra, generated this research
attempt on 2026-10-11. The exact serving revision is not exposed. Codex GPT-6
performed orchestration and repository-format review; its exact serving revision
is also not exposed. The mathematical route, boundary checks and final adversarial
self-audit followed contribution prompts A then B in
\texttt{docs/contribution\_prompts.md}, after reading \texttt{CONTRIBUTING.md}.
Python 3.14.2 assembled the standalone documents; any mathematical computation is
specified in Verification. No theorem prover was used. No human mathematical
authorship, substantial human proof correction, or independent human proof checking
is asserted. Model-assisted review, source retrieval, and structural checks do not
establish correctness or novelty. Full private reasoning transcripts are unavailable
as reproducibility artifacts; the complete arguments, inputs and outcomes actually
used are stated here or in the companion batch audit.
\section{Final status}
Unresolved. This is a research attempt with an explicit failed route or restricted
calculation, not a full solution, not a current-status certification, and not a claim
that the elementary material is new. The catalogue's original source-status label
is preserved independently of this attempt status.
\end{document}
